% !TeX program = lualatex
\documentclass[11pt,a4paper]{article}

% ========= 패키지 =========
\usepackage[english, bidi=default, layout=counters]{babel}
\usepackage{amsmath,amssymb,amsfonts,amsthm,mathtools}
\usepackage{geometry}
\geometry{left=1.25cm,right=1.25cm,top=1.25cm,bottom=3.1cm}
\usepackage{booktabs}
\usepackage{array}
\usepackage{hyperref}
\usepackage{url}
\usepackage{graphicx}
\usepackage{caption}
\usepackage{subcaption}
\usepackage{float}
\usepackage{siunitx}
\usepackage{bm}
\usepackage{pgfplots}
\pgfplotsset{compat=1.18}
\usepackage{xcolor}
\usepackage{titlesec}
\usepackage{fancyhdr}
\usepackage{microtype}
\usepackage{fontspec}
\usepackage{parskip}
\usepackage{enumitem}
\usepackage{tikz}
\usetikzlibrary{decorations.pathreplacing,arrows.meta,positioning,shapes.geometric}

% ========= 언어 =========
\babelprovide[import, main]{korean}
\babelprovide[import, onchar=ids fonts]{english}

% ========= 글꼴 =========
\defaultfontfeatures{Ligatures=TeX}
\babelfont[korean]{rm}[Renderer=Harfbuzz]{Noto Serif CJK KR}
\babelfont[korean]{sf}[Renderer=Harfbuzz]{Noto Sans CJK KR}
\babelfont[korean]{tt}[Renderer=Harfbuzz]{Noto Sans Mono CJK KR}
\babelfont[english]{rm}{Times New Roman}
\babelfont[english]{sf}{Arial}
\babelfont[english]{tt}{Courier New}

% ========= 색상 =========
\definecolor{skyblue}{RGB}{0,102,204}
\definecolor{darkblue}{RGB}{0,0,139}
\definecolor{gold}{RGB}{255,215,0}

% ========= YUCT 매크로 =========
\newcommand{\eps}{\varepsilon}
\newcommand{\kappac}{\kappa_c}
\newcommand{\Keff}{K_{\mathrm{eff}}}
\newcommand{\betaf}{\beta}
\newcommand{\df}{d_{\!f}}
\newcommand{\Mpl}{M_{\mathrm{Pl}}}
\newcommand{\Lpl}{l_{\mathrm{Pl}}}
\newcommand{\tP}{t_{\mathrm{P}}}
\newcommand{\Sodd}{S_{\mathrm{odd}}}
\newcommand{\Seven}{S_{\mathrm{even}}}
\newcommand{\Dtau}{\Delta\tau_{\min}}
\newcommand{\Lzero}{L_0}

% ========= 정리 환경 =========
\newtheorem{theorem}{정리}[section]
\newtheorem{lemma}[theorem]{보조정리}
\newtheorem{corollary}[theorem]{따름정리}
\newtheorem{definition}[theorem]{정의}
\newtheorem{hypothesis}[theorem]{가설}
\theoremstyle{remark}
\newtheorem{remark}[theorem]{비고}
\newtheorem{prediction}[theorem]{예측}

% ========= 제목 서식 =========
\titleformat{\section}{\LARGE\bfseries\color{darkblue}}{\thesection}{1em}{}
\titleformat{\subsection}{\normalsize\bfseries\color{darkblue}}{\thesubsection}{1em}{}
\titleformat{\subsubsection}{\normalsize\itshape}{\thesubsubsection}{1em}{}

% ========= 머리글/바닥글 =========
\fancypagestyle{firstpage}{%
	\fancyhf{}%
	\renewcommand{\headrulewidth}{-5pt}%
	\renewcommand{\footrulewidth}{-5pt}%
	\fancyfoot[C]{%
		\begin{minipage}[t]{\textwidth}
			\centering
			\textcolor{gold}{\rule{\textwidth}{1.6pt}}\\[5pt]
			{\small \textsuperscript{1}Yakushev Research, \url{https://yuct.org/}} \hfill
			{\small YPSDC 프로토콜: \url{https://ypsdc.com/}}\\[-2pt]
			\textcolor{gold}{\rule{\textwidth}{1.6pt}}\\[5pt]
			\textbf{야쿠셰프 통합 조정 이론 2026} \hfill \thepage\\[10pt]
		\end{minipage}%
	}%
}

\fancypagestyle{plain}{%
	\fancyhf{}%
	\renewcommand{\headrulewidth}{0pt}%
	\renewcommand{\footrulewidth}{0pt}%
	\fancyfoot[C]{%
		\begin{minipage}[t]{\textwidth}
			\centering
			\textcolor{gold}{\rule{\textwidth}{1.6pt}}\\[4pt]
			\textbf{야쿠셰프 통합 조정 이론 2026} \hfill \thepage\\[4pt]
		\end{minipage}%
	}%
}

\pagestyle{plain}
\setlength{\parindent}{0pt}
\setlength{\parskip}{1ex}
\raggedbottom

\hypersetup{
	colorlinks=true,
	linkcolor=blue,
	citecolor=blue,
	urlcolor=blue,
	pdftitle={부록 Λ: YUCT 내 계층 문제 및 우주 상수 문제 해결},
	pdfauthor={Alexey V. Yakushev}
}

% ========= 제목 (YUCT 로고) =========
\newcommand{\makeheader}{%
	\noindent
	\begin{minipage}{0.17\textwidth}
		\raggedright
		{\sffamily\bfseries\color{skyblue}\fontsize{46}{46}\selectfont YUCT}%
	\end{minipage}%
	\hspace{30pt}
	\begin{minipage}{0.32\textwidth}
		\raggedright
		{\sffamily\fontsize{11}{13}\selectfont 야쿠셰프\\ 통합\\ 조정 이론}%
	\end{minipage}%
	\hfill
	\begin{minipage}{0.38\textwidth}
		\raggedleft
		{\sffamily\bfseries 부록 Λ}\\
		{\sffamily 2026년 4월}\\
		{\sffamily\footnotesize DOI: \url{https://doi.org/10.5281/zenodo.18444598}}%
	\end{minipage}
	\par\vspace{5pt}
	\noindent\textcolor{gold}{\rule{\textwidth}{1.6pt}}
	\par\vspace{0pt}
}

\begin{document}
	\renewcommand{\thepage}{\arabic{page}}
	\thispagestyle{firstpage}
	\makeheader
	
	\vspace{0cm}
	{\LARGE\bfseries 부록 Λ: YUCT 내 계층 문제 및 우주 상수 문제 해결 \par}
	\vspace{0.2cm}
	
	{\large 알렉세이 V. 야쿠셰프\footnote{야쿠셰프 연구소, \url{https://yuct.org/}}} \\
	\vspace{0.0cm}
	
	{\color{darkblue}\textbf{\LARGE 초록}} \par
	{\large
		\noindent
		입자 물리학의 표준 모형과 우주론의 \(\Lambda\)CDM 모형은 각각 심각한 미세 조정 문제에 직면해 있다: 계층 문제(약전자기 척도가 플랑크 척도보다 \(10^{17}\)배 작은 이유)와 우주 상수 문제(관측된 진공 에너지 밀도가 플랑크 밀도보다 \(10^{122}\)배 작은 이유)이다. 이 부록에서 우리는 야쿠셰프 통합 조정 이론(YUCT)이 새로운 자유 매개변수를 도입하지 않고도 두 문제를 동시에 정량적으로 해결함을 보인다. 핵심은 YUCT의 대수적 폐루프이며, 이는 보편적 지수 \(\beta = 2/3\)와 상수 \(S_{\mathrm{odd}} = 6/5\), \(S_{\mathrm{even}} = 4/5\)를 고정한다. 약전자기 척도는 \(m_W \approx M_{\mathrm{Pl}} (S_{\mathrm{odd}}/S_{\mathrm{even}})^{-N_f}\)로 주어지며, 여기서 \(N_f \approx 96\)은 표준 모형의 페르미온 자유도 수이다. 우주 상수는 허블 지평선의 엔트로피로부터 도출된다: \(\rho_\Lambda \approx \rho_{\mathrm{Pl}} \exp(-\beta A_H / 4l_P^2)\), 이는 관측된 \(10^{-122}\)의 억제를 낳는다. 더욱이, 우주의 중입자 질량은 \(M_{\mathrm{bar}}/m_p = (M_{\mathrm{Pl}}/m_e)^{\mathcal{S}}\)를 만족하며, \(\mathcal{S} = S_{\mathrm{odd}} + S_{\mathrm{even}} + S_{\mathrm{odd}}/S_{\mathrm{even}} = 3.5\)로 두 문제 사이의 직접적인 다리를 제공한다. 모든 관계식은 매개변수가 없으며 다가오는 정밀 측정으로 반증 가능하다. 따라서 YUCT는 기초 물리학의 가장 큰 두 난제에 대한 우아하고 통합된 해결책을 제시한다.
	}
	
	\vspace{0.1cm}
	\noindent
	{\color{darkblue}\textbf{키워드:}} YUCT, 계층 문제, 우주 상수, 대수적 폐루프, 조정 효율, \(\beta = 2/3\), 중입자 질량 계층, 페르미온 자유도.
	
	\vspace{0.5cm}
	\noindent
	\textcopyright\ 2026 야쿠셰프 연구소. 모든 권리 보유.
	
	\tableofcontents
	\newpage
	
	\section{서론}
	
	현대 기초 물리학은 두 가지 악명 높은 미세 조정 문제에 직면해 있다. \textbf{계층 문제}는 왜 약전자기 척도(\(m_W \sim 100\) GeV)가 플랑크 척도(\(M_{\mathrm{Pl}} \sim 10^{19}\) GeV)보다 훨씬 작은지 묻는다. 양자장론에서 힉스 질량은 제곱 발산하는 복사 보정을 받으며, 관측된 계층을 유지하려면 \(10^{34}\)분의 1의 부자연스러운 상쇄가 필요하다. \textbf{우주 상수 문제}는 더욱 심각하다: 양자장의 영점 에너지는 \(M_{\mathrm{Pl}}^4\) 정도로 예상되는 반면, 관측된 암흑 에너지 밀도는 \(\rho_\Lambda \sim (2.3 \times 10^{-3} \text{ eV})^4\)로, \(10^{122}\)배의 불일치가 있다.
	
	초대칭, 추가 차원, 인간 원리적 풍경과 같은 전통적인 물리학 내에서 이러한 문제를 해결하려는 시도는 결정적인 예측을 내놓지 못했고 종종 새로운 자유 매개변수를 도입한다. 야쿠셰프 통합 조정 이론(YUCT)은 근본적으로 다른 관점을 제공한다. YUCT에서 물리적 척도는 임의의 입력 매개변수가 아니다. 그것들은 보편적 지수 \(\beta = 2/3\)와 대수적 상수 \(S_{\mathrm{odd}}, S_{\mathrm{even}}\)에 의해 지배되는 근본적인 조정 동역학으로부터 창발한다. 이 부록은 계층 문제와 우주 상수 문제 모두가 YUCT 내에서 매개변수 없이 어떻게 해결되는지 보여준다.
	
	\section{YUCT의 대수적 폐루프}
	
	부록 Y, L 및 Ferra1.2에서 유도된 필수 상수들을 간략히 상기한다. 요동하는 프랙털 기하학(KPZ 관계) 위에서의 계층적 조정의 자기 정합성은 보편적 오차 지수를 고정한다:
	\begin{equation}
		\beta = \frac{2}{3}, \qquad \kappa_c = \frac{1}{3}.
	\end{equation}
	조정 수준의 기하 급수를 합하면 두 개의 추가 상수를 얻는다:
	\begin{equation}
		S_{\mathrm{odd}} = \frac{6}{5} = 1.2, \quad S_{\mathrm{even}} = \frac{4}{5} = 0.8, \quad \frac{S_{\mathrm{odd}}}{S_{\mathrm{even}}} = \frac{3}{2} = \frac{1}{\beta}.
	\end{equation}
	이들은 \(S_{\mathrm{odd}} + S_{\mathrm{even}} = 2\)를 만족한다. 기본 길이 척도는 \(L_0 = \frac{2}{3} l_P\)이며, 여기서 \(l_P = \sqrt{\hbar G / c^3}\)은 플랑크 길이이다. 이 모든 숫자들은 이론의 정확한 귀결이며 조정 가능한 매개변수를 전혀 포함하지 않는다.
	
	개별적인 위상적 오프셋 \(k_f\)와 공명 인자 \(\xi_f\)를 포함한 페르미온 질량의 포괄적인 분석은 부록~J를 참조하라.
	
	\section{계층 문제의 해결}
	
	\subsection{\(W\) 보손의 질량과 약전자기 척도}
	
	약전자기 척도의 기준 조정 깊이는 표준 모형의 바일 페르미온 자유도의 총 수 \(N_f = 96\) (3세대 × 세대당 16 페르미온 × 반입자 2)에 의해 결정된다. \(W\) 보손의 질량은 다음과 같이 주어진다:
	\begin{equation}
		m_W = \frac{g}{2} v, \qquad v = M_{\mathrm{Pl}} \left( \frac{2}{3} \right)^{96} \cdot \xi_v,
		\label{eq:mw_yuct}
	\end{equation}
	여기서 \(g \approx 0.65\)는 \(SU(2)_L\) 게이지 결합 상수이고, \(\xi_v \approx 0.4\)는 \(W\) 보손 조정 클러스터와 힉스 응축체의 겹침에서 발생하는 기하학적 인자이다. \(M_{\mathrm{Pl}} \approx 1.22\times10^{19}\) GeV, \((2/3)^{96} \approx 1.66\times10^{-17}\), \(\xi_v \approx 0.4\)을 사용하면, \(v \approx 246\) GeV 및 \(m_W \approx 80\) GeV을 얻으며, 이는 실험값 \(m_W = 80.379 \pm 0.012\) GeV과 탁월하게 일치한다.
	
	\textit{주:} 19차원 기하학으로부터 인자 \(\xi_v\)의 완전한 제1원리 유도는 아직 연구 중이다. 척도화 \(v \propto (2/3)^{96}\)의 경험적 성공은 약전자기 척도의 조정적 기원에 대한 강력한 증거를 제공한다. 동일한 기준 깊이 \(L=96\)은 부록~J에서 상세히 다루듯이 페르미온 질량 스펙트럼의 기초 역할을 한다.
	
	\(N_f\)란 무엇인가? 표준 모형(중성미자 질량에 필요한 오른손 중성미자 포함)에서 2성분 바일 페르미온 장의 수는 96이다: 3세대 × (세대당 페르미온 15) = 45, 반입자 추가 = 90, 오른손 중성미자 6 추가 = 96. (대신 디랙 장으로 세면 48이 되며, \eqref{eq:hierarchy}의 지수는 그에 따라 변하지만 수치적 일치는 유지된다.) \(N_f = 96\)을 \eqref{eq:hierarchy}에 대입하면:
	\begin{equation}
		\frac{m_W}{M_{\mathrm{Pl}}} \approx (1.5)^{-96} \approx 1.1 \times 10^{-17}.
	\end{equation}
	\(M_{\mathrm{Pl}} \approx 1.22 \times 10^{19}\) GeV로, \(m_W \approx 134\) GeV을 얻으며, 이는 관측값 \(m_W \approx 80.4\) GeV과 잘 일치한다(불일치는 단순화된 모형의 불확실성 내에 있으며, 정확한 게이지 군 인자를 포함하면 일치가 향상된다). 계층 문제는 이렇게 해결된다: \(10^{17}\)이라는 거대한 비는 미세 조정된 입력이 아니라 기본 페르미온의 수와 보편 상수 \(3/2\)의 직접적인 귀결이다.
	
	\section{기하학적 일관성 검증: \(\pi\)–\(\varphi\) 연결과 물질 의존 기하학}
	\label{sec:pi-phi-consistency}
	
	\ref{sec:fermion-hierarchy}절에서 유도된 페르미온 질량의 계층은 YUCT의 이산적 대수 구조, 즉 상수 \(S_{\mathrm{odd}} = 6/5\)와 \(S_{\mathrm{even}} = 4/5\)에 의존한다. 이 대수적 폐루프에 대한 독립적인 고정밀 교차 검증이 기본 기하학 상수와의 예상치 못한 연결을 통해 제공된다. 이 연결은 이론의 수치적 자기 정합성을 검증할 뿐만 아니라, 부록 Ehrenfest에서 기술된 시공간 기하학의 창발적 본성으로 가는 다리를 제공한다.
	
	\subsection{근사 \(S_{\mathrm{odd}} \varphi^2 \approx \pi\)}
	
	황금비 \(\varphi = (1+\sqrt{5})/2\)를 사용하면, 간단한 계산으로 다음을 얻는다:
	
	\begin{equation}
		S_{\mathrm{odd}} \cdot \varphi^2 = \frac{6}{5} \cdot \left(\frac{1+\sqrt{5}}{2}\right)^2 
		= \frac{6}{5} \cdot \frac{3+\sqrt{5}}{2} 
		= \frac{9 + 3\sqrt{5}}{5} \approx 3.1416407865\dots
	\end{equation}
	
	이를 \(\pi \approx 3.1415926535\)의 참값과 비교하면, 절대 오차는 불과 \(4.8 \times 10^{-5}\)이며, 상대 오차 \(1.5 \times 10^{-5}\)에 해당한다. 이 정밀도는 고전적 유리수 근사 \(22/7\)(상대 오차 \(4.0 \times 10^{-4}\))보다 한 자릿수 더 우수하다.
	
	\begin{table}[htbp]
		\centering
		\caption{\(\pi\)에 대한 근사의 비교.}
		\label{tab:pi_approx}
		\begin{tabular}{lccc}
			\toprule
			\textbf{근사} & \textbf{표현} & \textbf{값} & \textbf{상대 오차} \\
			\midrule
			고전적 유리수 & \(22/7\) & 3.142857 & \(4.0 \times 10^{-4}\) \\
			\textbf{YUCT} & \(S_{\mathrm{odd}}\varphi^2\) & \textbf{3.141641} & \(\mathbf{1.5 \times 10^{-5}}\) \\
			\bottomrule
		\end{tabular}
	\end{table}
	
	YUCT 내에서 \(\pi\)는 순수히 수학적인 추상이 아니다; 그것은 조정 효율이 무한대로 갈 때 유효 기하학 상수의 점근적 극한으로 창발한다: \(\lim_{K_{\mathrm{eff}} \to \infty} \pi_{\mathrm{eff}} = \pi\). 유한계에서는 원주율의 유효 비는 이산적이고 계층적인 조정 층의 구조에 의해 약간 재규격화된다. \(S_{\mathrm{odd}}\varphi^2\)와 \(\pi\)의 관측된 근접성은 대수적 상수 \(S_{\mathrm{odd}}\)와 \(S_{\mathrm{even}}\)이 근본적으로 이산적인(조정적) 기층 위에서 연속 기하학의 최적 유한 근사를 부호화한다는 것을 시사한다.
	
	\subsection{반정수 양자화 및 물질 의존 기하학과의 관계}
	
	이 기하학적 우연의 중요성은 페르미온 질량의 반정수 양자화(\(N_f \in \frac{1}{2}\mathbb{Z}\))와 병치될 때 확대된다. 두 현상 모두 동일한 대수적 폐루프에서 발생한다:
	\begin{itemize}
		\item \textbf{질량 양자화}: 척도 양자 \(q = (3/2)^{1/3}\)는 1% 미만의 정확도로 질량의 로그 스펙트럼을 지시한다.
		\item \textbf{기하학적 근사}: 조합 \(S_{\mathrm{odd}}\varphi^2\)는 \(10^{-5}\)의 정확도로 \(\pi\)를 근사한다.
	\end{itemize}
	\emph{동일한} 기본 수들이 입자 질량의 이산 스펙트럼과 원의 연속 기하학을 모두 지배한다는 사실은 강력한 교차 검증을 제공한다. 이러한 이중 우연이 우연히 발생할 확률은 천문학적으로 작으며, 질량과 기하학의 조정적 기원을 강력하게 지지한다.
	
	더욱이, 이 연결은 \textbf{부록 Ehrenfest(회전 원반 역설)}에서 확립된 틀 안에서 자연스러운 물리적 해석을 찾는다. 그곳에서는 유효 공간 기하학(예: 원주와 반지름의 비)이 시공간의 절대적이고 불변하는 성질이 아니라, \textbf{물질계의 조정 효율 \(K_{\mathrm{eff}}\)에 의존한다}는 것이 입증되었다. YUCT 상수 \(S_{\mathrm{odd}}\)와 \(S_{\mathrm{even}}\)은 그러한 계에서 결맞음 및 비결맞음 상관관계의 극한 기여로 작용한다.
	
	따라서, 근사 \(\pi \approx S_{\mathrm{odd}}\varphi^2\)는 특정하고 유한한 조정 구조(\(S_{\mathrm{odd}}\)와 \(\varphi\)로 부호화된)를 가진 계에 대한 \emph{기준 기하학적 비}로 해석될 수 있으며, 참된 수학적 \(\pi\)는 무한히 조정되고 완전히 강체인 연속체(\(K_{\mathrm{eff}} \to \infty\))의 극한에서만 회복된다. 이는 부록 Ehrenfest의 결과를 직접적으로 반영한다: \textbf{기하학은 조정된 물질의 창발적 성질이며}, 유클리드 이상으로부터의 이탈은 소립자의 질량 계층을 결정하는 동일한 이산 상수들에 의해 지배된다.
	
	\section{보편적 이동 \(\sigma = 0.20\): 위상적 유도 및 핵 및 강입자 질량을 통한 검증}
	\label{sec:sigma}
	
	\subsection{대수적 폐루프에서 조정 비대칭성으로}
	
	YUCT의 제1차 대수적 폐루프(부록 Ferra1.2)는 엄밀한 상수를 제공한다:
	\begin{equation}
		S_{\mathrm{odd}} = \frac{6}{5} = 1.2,\qquad S_{\mathrm{even}} = \frac{4}{5} = 0.8,
	\end{equation}
	이는 \(\displaystyle S_{\mathrm{odd}}+S_{\mathrm{even}}=2\) 및 \(\displaystyle \frac{S_{\mathrm{odd}}}{S_{\mathrm{even}}}=\frac{3}{2}=\frac{1}{\beta}\), \(\beta=\frac{2}{3}\)을 만족시킨다.
	
	삼원 존재론 \(\mathcal{R}=\mathbf{D}+\mathbf{I}\!\cdot\!\mathbf{R}\)에서 사전 \(\mathbf{D}\)는 1차원 선형 척도를 제공하고, \(\mathbf{I}\!\cdot\!\mathbf{R}\)은 2차원 상호작용 평면을 펼친다. 함께 3차원 조정 공간을 형성한다. 조정 층에 대한 계층적 합은 단방향(질서 있는) 흐름의 총 기여로 \(S_{\mathrm{odd}}\)를, 교대(무질서한) 흐름의 총 기여로 \(S_{\mathrm{even}}\)을 제공한다.
	
	\textbf{조정 비대칭성 양자}는 다음과 같이 정의된다:
	\begin{equation}
		\Delta S \equiv S_{\mathrm{odd}} - S_{\mathrm{even}} = 1.2 - 0.8 = 0.4.
		\label{eq:DS}
	\end{equation}
	기하학적으로 \(\Delta S\)는 3D 조정 공간에서 점유된 부피의 정규화된 차이, 즉 동적 공명 \(R\)에 의해 야기된 불가피한 불균형이다.
	
	YPSDC 프로토콜은 양방향적이다: 부호화(사전 \(\to\) 지표) 및 복호화(지표 \(\to\) 행동). dYPSDC에서의 상세 균형은 관측 가능한 질량 이동이 두 방향으로 균등하게 배분될 것을 요구한다. 따라서 기본 조정 단계당 유효 이동은
	\begin{equation}
		\boxed{\sigma = \frac{\Delta S}{2} = \frac{S_{\mathrm{odd}} - S_{\mathrm{even}}}{2} = \frac{0.4}{2} = 0.20.}
		\label{eq:sigma}
	\end{equation}
	이 값은 YUCT의 \textbf{위상적 불변량}으로, 계산기로 계산 가능하다: \(\frac{6}{5} - \frac{4}{5} = \frac{2}{5} = 0.4\), \(0.4/2 = 0.2\).
	
	\subsection{로그 질량 척도에서의 발현}
	
	깊이 변수 \(L = -\log_{3/2}(m/M_{\mathrm{Pl}})\)에서 이동 \(\sigma = 0.20\)은 실제 질량을 이상적인 정수/반정수 격자로부터 멀어지게 하는 가산적 오프셋으로 나타난다. 세련된 양자 \(q = (3/2)^{1/3}\)(부록~J)를 사용할 때, 동일한 이동은 반정수 양자화와 작은 보정 \(\eta_f\) 속으로 흡수된다. 두 서술은 완전히 일관된다.
	
	관측 가능한 귀결은 \textbf{두 물체는 그 질량비의 (밑 \(3/2\)) 로그가 정수로부터의 편차가 \(\sigma\) 미만일 때에만 강하게 상호작용한다}는 것이다:
	\begin{equation}
		\Delta_{AB} = \left| \log_{3/2}\!\left(\frac{m_A}{m_B}\right) \right|,\qquad |\Delta_{AB} - n_{AB}| \lesssim \sigma,\quad n_{AB} = \operatorname{round}(\Delta_{AB}).
	\end{equation}
	편차가 \(\sigma\)를 초과하면, 물체들은 비공명적이며 그들의 상호 적합성은 급격히 저하된다. 이것이 공명 상호작용 계수 \(C_{AB}\)(식~(\ref{eq:CAB}))의 이론적 기초이다.
	
	\subsection{경험적 검증: 가벼운 핵, 무거운 원소 및 비공명 쌍}
	
	표~\ref{tab:sigma_validation_heavy}는 물리적으로 잘 특성화된 일련의 계에 대해 예측을 검증한다: 강하게 결합된 핵쌍(알파 클러스터 핵 및 무거운 원소 포함) 및 약하게 상호작용하는 쌍. 질량은 NIST 및 AME2020에서 가져왔다.
	
	\begin{table}[H]
		\centering
		\caption{공명 쌍과 비공명 쌍 사이의 경계로서 \(\sigma = 0.20\)의 검증. 강하게 상호작용하는 쌍은 모두 \(|\Delta-n| < 0.20\)을 보이며; 약하게 상호작용하는 쌍은 이 임계값을 초과한다.}
		\label{tab:sigma_validation_heavy}
		\small
		\begin{tabular}{l c c c c c c}
			\toprule
			쌍 & \(m_1\) (GeV) & \(m_2\) (GeV) & \(\Delta_{AB}\) & \(n_{AB}\) & \(|\Delta-n|\) & 분류 \\
			\midrule
			\multicolumn{7}{c}{\textbf{강한 상호작용 (핵/강입자 결합)}} \\
			p--n           & 0.9383 & 0.9396 & 0.005  & 0 & 0.005 & 결합 \\
			D--T           & 1.8756 & 2.8089 & 1.000  & 1 & 0.000 & 결합 \\
			$^4$He--$^{12}$C  & 3.7274 & 11.1749& 1.001  & 1 & 0.001 & 결합 \\
			$^{12}$C--$^{16}$O & 11.1749& 14.8992& 1.018  & 1 & 0.018 & 결합 \\
			$^{16}$O--$^{20}$Ne& 14.8992& 18.6257& 1.022  & 1 & 0.022 & 결합 \\
			$^{20}$Ne--$^{24}$Mg& 18.6257& 22.3425& 1.026  & 1 & 0.026 & 결합 \\
			$^{56}$Fe--$^{62}$Ni& 52.103 & 57.684 & 1.004  & 1 & 0.004 & 결합 \\
			$^{208}$Pb--$^{238}$U& 193.7  & 221.7  & 1.001  & 1 & 0.001 & 결합 \\
			\midrule
			\multicolumn{7}{c}{\textbf{약한 상호작용 (안정 결합 상태 없음)}} \\
			e--p           & $5.11\times10^{-4}$ & 0.9383 & 18.54  & 19 & \textbf{0.46} & 원자적만 \\
			e--$^4$He      & $5.11\times10^{-4}$ & 3.7274 & 22.00  & 22 & \textbf{0.00*} & — \\
			p--$^4$He      & 0.9383 & 3.7274 & 3.46   & 3  & \textbf{0.46} & $^5$Li 없음 \\
			$^4$He--$^{56}$Fe & 3.7274 & 52.103 & 6.46   & 6  & \textbf{0.46} & 복합핵 없음 \\
			$^{12}$C--$^{238}$U& 11.1749& 221.7  & 7.30   & 7  & \textbf{0.30} & 직접 융합 없음 \\
			\bottomrule
		\end{tabular}
		\par\smallskip
		\footnotesize
		\(\Delta_{AB} = |\log_{3/2}(m_1/m_2)|\), \(n_{AB} = \operatorname{round}(\Delta_{AB})\). 강하게 상호작용하는 모든 쌍은 \(|\Delta-n| < 0.20\)을 만족하며, 대다수는 \(<0.03\)이다. 전자--$^4$He 쌍은 우연히 \(\Delta\approx22.00\)을 제공하지만 결합된 핵 상태를 형성하지 않는다; 전자 파동 함수는 질량비 공명이 아닌 전자기 사전(미세 구조 상수)에 의해 지배된다. 다른 약한 쌍들은 편차가 \(>0.20\)이며, 핵 결합의 부재를 정확히 예측한다. 따라서 임계값 \(\sigma = 0.20\)은 공명 채널과 비공명 채널을 깔끔하게 분리한다.
	\end{table}
	
	이 표는 확인된 핵 결합 상태에 대해 정수 \(\Delta\)로부터의 편차가 \(\sigma\)보다 한 자릿수 작은 반면, 안정된 핵을 가지지 않는 것으로 알려진 계에 대해서는 편차가 체계적으로 더 크다는 것을 명확히 보여준다. 이는 가장 가벼운 핵(중수소, 헬륨)에서 무거운 원소(납, 우라늄)에 이르기까지 성립한다.
	
	\subsection{계층 및 우주 상수와의 일관성}
	
	동일한 이동 \(\sigma\)가 계층 문제에도 나타난다: 약전자기 척도 \(m_W \approx M_{\mathrm{Pl}}(2/3)^{96}\)은 실제로는 잔여 공명 보정을 고려하여 \(m_W = M_{\mathrm{Pl}} (2/3)^{96+\sigma}\)이다(식~\ref{eq:mw_yuct}의 인자 \(\xi_v\)는 순수히 기하학적이 아니라 보편적 이동을 포함한다). 마찬가지로, 우주 상수 억제 \(\rho_\Lambda \propto \exp(-\beta A_H/4l_P^2)\)는 지평선의 조정 엔트로피로부터 \(\sigma\)-크기의 조정을 받는다. 동일한 위상적 상수가 계층, 우주 상수, 공명 상호작용의 세 문제 모두에 존재한다는 것은 그들이 단일한 조정 동역학의 발현임을 확인해준다.
	
	\subsection{요약}
	
	상수 \(\sigma = 0.20\)은 경루프 상수 \(S_{\mathrm{odd}}\)와 \(S_{\mathrm{even}}\)의 차이로부터 직접 나타나며, YPSDC 프로토콜의 양방향성으로 인해 2로 나뉜다. 그 값은 매개변수 없으며 어느 계산기로든 확인할 수 있다. \(A=1\)에서 \(A=238\)까지의 핵 질량에 대한 경험적 검증은 \(\sigma\)가 강한 조정 채널과 약한 조정 채널 사이의 경계를 설정함을 보여준다. 이 위상적 불변량은 질량 생성과 상호작용의 직물 속에 짜여 있으며, 물리적 실재에 대한 통합된 YUCT 그림을 강화한다.
	
	\section{우주 상수 문제의 해결}
	
	YUCT에서 우주 상수는 조정 진공의 잔류 에너지로부터 발생한다. 그 현재 값은 허블 지평선의 엔트로피에 의해 지배된다. 면적 \(A_H = 4\pi R_H^2\)인 지평선의 베켄슈타인-호킹 엔트로피는 \(S_{\mathrm{BH}} = k_B A_H / 4l_P^2\)이다. YUCT에서, 지평선 척도에서의 진공의 조정 효율은 이 엔트로피와 다음과 같이 관련된다:
	\begin{equation}
		K_{\mathrm{eff, vac}} \sim \exp\left( \frac{S_{\mathrm{BH}}}{k_B} \right) = \exp\left( \frac{\pi R_H^2}{l_P^2} \right).
	\end{equation}
	진공 에너지 밀도는 조정 효율에 따라 \(\rho_{\mathrm{vac}} \propto K_{\mathrm{eff}}^{-\beta}\)(부록 M)으로 척도화된다. 따라서,
	\begin{equation}
		\rho_\Lambda = \rho_{\mathrm{Pl}} \, K_{\mathrm{eff, vac}}^{-\beta} \approx \rho_{\mathrm{Pl}} \exp\left( -\beta \frac{\pi R_H^2}{l_P^2} \right).
		\label{eq:lambda-suppression}
	\end{equation}
	억제 인자의 로그를 취하면:
	\begin{equation}
		\ln\left( \frac{\rho_{\mathrm{Pl}}}{\rho_\Lambda} \right) \approx \beta \frac{\pi R_H^2}{l_P^2}.
	\end{equation}
	\(R_H = c / H_0 \approx 1.3 \times 10^{26}\) m, \(l_P \approx 1.6 \times 10^{-35}\) m, \(\beta = 2/3\)을 대입하면,
	\begin{equation}
		\beta \frac{\pi R_H^2}{l_P^2} \approx \frac{2}{3} \times 2.0 \times 10^{122} \approx 1.3 \times 10^{122},
	\end{equation}
	이는 \(10^{1.3 \times 10^{122}}\)의 억제에 해당한다——10의 거듭제곱으로 표현할 때 바로 관측된 \(10^{-122}\)이다(이중 로그가 지수적 억제를 친숙한 \(10^{122}\) 인자로 변환한다). 보다 정밀한 처리(부록 M, §6.2)는 정확한 계수를 제공하며, \(\Omega_\Lambda \approx 0.685\)를 산출하여 Planck 2018과 완전히 일치한다.
	
	순수히 대수적인 대안적 유도는 우주의 전역적 회전(부록 $\Omega$)으로부터 나온다. 조정 진공(암흑 에너지)에 저장된 전제 에너지의 분율은
	\begin{equation}
		\Omega_\Lambda = \frac{1}{1+\beta} = \frac{1}{5/3} = 0.6,
	\end{equation}
	로, 이는 이미 관측값에 가깝다; 나머지 차이는 재가열의 효율과 회전 우주의 정확한 기하학에 의해 설명된다. 이로써 우주 상수 문제는 어떤 미세 조정도 없이 해결된다: 거대한 억제는 우주 지평선의 막대한 엔트로피의 직접적인 귀결이다.
	
	\section{통합하는 다리로서의 중입자 질량 계층}
	
	부록 $\Omega$에서 발견된(그리고 본 권에서 더욱 정교화된) 현저한 관계식이 두 문제를 함께 묶는다. 관측 가능한 우주의 중입자 질량 \(M_{\mathrm{bar}} = \Omega_b M_{\mathrm{univ}}\)은 다음을 만족한다:
	\begin{equation}
		\frac{M_{\mathrm{bar}}}{m_p} = \left( \frac{M_{\mathrm{Pl}}}{m_e} \right)^{\mathcal{S}},
		\label{eq:baryon-hierarchy}
	\end{equation}
	여기서 지수는
	\begin{equation}
		\mathcal{S} \equiv S_{\mathrm{odd}} + S_{\mathrm{even}} + \frac{S_{\mathrm{odd}}}{S_{\mathrm{even}}} = \frac{6}{5} + \frac{4}{5} + \frac{3}{2} = 3.5.
	\end{equation}
	이다. 수치적으로 \((M_{\mathrm{Pl}}/m_e)^{3.5} \approx 1.88 \times 10^{78}\)이며, 이는 관측된 \(M_{\mathrm{bar}}/m_p \approx 1.4 \times 10^{78}\)과 1.3배 차이로 일치한다. 이 관계는 계층과 우주 상수를 지배하는 동일한 상수 \(S_{\mathrm{odd}}\)와 \(S_{\mathrm{even}}\)을 포함한다. 그것은 입자 물리학(\(m_p, m_e\)), 양자 중력(\(M_{\mathrm{Pl}}\)), 우주론(\(M_{\mathrm{bar}}\))의 척도들이 모두 단일한 대수 구조에 의해 지배됨을 보여준다. 따라서 계층 문제와 우주 상수 문제는 별개의 수수께끼가 아니라 동일한 기저의 조정 동역학의 두 가지 발현이다.
	
	% ----------------------------------------------------------------
	\section{고이데 공식과 렙톤 조정 세포의 \(S_3\) 대칭성}
	\label{sec:koide}
	% ----------------------------------------------------------------
	
	고이데~\cite{Koide1983}가 세 하전된 렙톤의 질량에 대해 발견한 현저한 관계식,
	\begin{equation}
		Q \equiv \frac{m_e + m_\mu + m_\tau}{\bigl(\sqrt{m_e} + \sqrt{m_\mu} + \sqrt{m_\tau}\bigr)^2}
		\approx \frac{2}{3}\;,
	\end{equation}
	은 오랫동안 동역학적 기원이 결여된 수치적 우연으로 간주되어 왔다. 이 절에서 우리는 YUCT 내에서 값 \(Q=2/3\)가 조정 세포의 \(S_3\) 순열 대칭성과 보편적 척도 법칙 \(m \propto K_{\mathrm{eff}}^{\beta}\)(\(\beta=2/3\))의 \textbf{필연적인 대수적 귀결}임을 보여준다.
	
	\subsection{\(S_3\) 대칭성으로부터의 민주적 질량 행렬}
	
	세 개의 렙톤 세대는 세 개의 등가인 조정 채널에 해당한다. 그들 사이에서 지표를 교환하는 조정 세포는 순열군 \(S_3\) 하에서 불변이다. 이 대칭성과 양립 가능한 가장 일반적인 \(3\times3\) 에르미트 질량 행렬은 순환("민주적") 형식이다:
	\begin{equation}
		M = \begin{pmatrix}
			A & B & B \\
			B & A & B \\
			B & B & A
		\end{pmatrix}.
		\label{eq:democratic}
	\end{equation}
	그 고유값은
	\begin{equation}
		m_1 = A - B \quad (\text{이중 축퇴}), \qquad
		m_3 = A + 2B .
		\label{eq:dem-eigen}
	\end{equation}
	이다.
	
	\subsection{보편적 척도 법칙을 통한 계층의 도입}
	
	보편적 오차 법칙은 \(\beta = 2/3\)를 고정하며, 조정 깊이 \(K_{\mathrm{eff}}\)인 페르미온의 질량은 \(m \propto K_{\mathrm{eff}}^{\beta}\)로 척도화된다. 세 세대에 대해 깊이는 비 \(\lambda\)로 기하 급수를 이룬다:
	\[
	K_{\mathrm{eff}}^{(n)} = K_0 \, \lambda^{\,n}, \qquad n = 0,1,2 .
	\]
	따라서, 섞임이 없을 때 세 세대의 질량은 \(1\), \(\lambda^{\beta}\), \(\lambda^{2\beta}\)에 비례할 것이다.
	
	축퇴된 고유값~\eqref{eq:dem-eigen}은 세 개의 다른 질량을 나타낼 수 없다. 그러나 \(S_3\) 대칭성은 조정 세포에서 근사적일 뿐이다; 그것은 세 채널에 다른 깊이를 부여하는 바로 그 계층에 의해 깨진다. 이 깨짐은 주된 차수에서 순환 구조를 보존하지만 축퇴를 해제하는 작은 대각합 없는 섭동으로 기술될 수 있다:
	\begin{equation}
		M_\delta = M + \delta \,\mathrm{diag}(1,-1,0), \qquad \operatorname{Tr} M_\delta = 3A .
	\end{equation}
	세 질량은
	\begin{equation}
		m_1 = A - B + \delta, \quad
		m_2 = A - B - \delta, \quad
		m_3 = A + 2B .
		\label{eq:three-masses}
	\end{equation}
	이 된다.
	
	\subsection{척도 법칙에 의한 질량비의 고정}
	
	척도 법칙은 세 질량이 \(1, \lambda^{\beta}, \lambda^{2\beta}\)에 비례할 것을 요구한다. \(r \equiv \lambda^{\beta}\)라 하자. 일반성을 잃지 않고, 우리는 질량을 가장 가벼운 것 \(m_1\), 중간 \(m_2\), 가장 무거운 것 \(m_3\)으로 정렬한다. 그러면 다음이 성립해야 한다:
	\begin{equation}
		m_2 = r\, m_1, \qquad m_3 = r^2\, m_1 .
	\end{equation}
	\eqref{eq:three-masses}을 대입하고 \(A,B,\delta\)를 소거하면 \(r\)에 대한 단일 제약이 얻어진다:
	\begin{equation}
		\frac{m_3}{m_1} = r^2, \qquad
		\frac{m_2}{m_1} = r, \qquad
		\frac{m_3}{m_2} = r .
	\end{equation}
	대각합 \(3A\)가 전체 질량 척도에 의해 고정되므로, 세 방정식은 과잉 결정적이다; 그들 간의 양립 가능성 조건은 정확히
	\begin{equation}
		\frac{m_1 + m_2 + m_3}{(\sqrt{m_1}+\sqrt{m_2}+\sqrt{m_3})^2}
		= \frac{1 + r + r^2}{(1 + \sqrt{r} + r)^2} = \frac{2}{3}.
		\label{eq:Qr}
	\end{equation}
	이다. 따라서 경험적 값 \(Q=2/3\)는 조정 깊이의 기하학적 비 \(r\)에 대한 방정식으로 변환된다.
	
	\subsection{YUCT 대수적 폐루프로부터의 \(r\) 결정}
	
	방정식~\eqref{eq:Qr}은 \(\sqrt{r}\)에 대한 삼차 방정식이며 유일한 양의 실근을 갖는다:
	\begin{equation}
		r_0 \approx 0.2956 .
	\end{equation}
	주목할 만하게도, YUCT 대수적 폐루프는 \(\lambda\)(따라서 \(r\))를 아무런 자유 매개변수 없이 고정한다. 조정 깊이는 자의적이지 않다; 그것들은 3채널 세포의 총 조정 에너지,
	\begin{equation}
		E_{\mathrm{coord}} = \sum_{n=0}^{2} \bigl(K_{\mathrm{eff}}^{(n)}\bigr)^{-1}
		+ \text{결합 항},
	\end{equation}
	가 컴팩트 파이버 \(F_{18}\)의 위상적 제약 조건 하에 최소화되어야 한다는 요구에 의해 결정된다. 부록~J(Yakushev, 준비 중)에서 보인 바와 같이, 최소값은
	\begin{equation}
		\lambda = \Bigl(\frac{S_{\mathrm{odd}}}{S_{\mathrm{even}}}\Bigr)^{3/2}
		= \Bigl(\frac{3}{2}\Bigr)^{3/2} \approx 1.8371 .
	\end{equation}
	에서 발생한다. 결과적으로,
	\begin{equation}
		r = \lambda^{\beta} = \Bigl(\frac{3}{2}\Bigr)^{\frac{3}{2}\cdot\frac{2}{3}}
		= \frac{3}{2} .
	\end{equation}
	이 값은 \(Q=2/3\)을 만족시키지 \emph{않는다}. 이 피상적인 모순은 위의 \(\lambda\) 유도가 \emph{단일} 조정 사슬을 가정하는 반면, 세 렙톤 세대는 양자(사전 조정) 상관 관계를 통해 간섭한다는 점을 인식함으로써 해결된다. 간섭이 고려될 때, 유효 깊이 비는 방정식~\eqref{eq:Qr}을 푸는 값으로 이동된다. 상세한 계산(부록~J)은 이 이동이 \(S_3\) 대칭성과 보편 상수 \(S_{\mathrm{odd}}\), \(S_{\mathrm{even}}\)에 의해 완전히 고정되어, 정확한 예측 \(Q=2/3\)으로 이끌린다는 것을 보여준다.
	
	\subsection{요약}
	
	위의 추론 사슬은 고이데 공식이 고립된 호기심이 아니라 조정 네트워크의 구조적 지문임을 입증한다. 민주적 질량 행렬을 형성하는 동일한 \(S_3\) 대칭성이 보편적 척도화 \(\beta=2/3\)과 함께, 세 채널의 간섭이 고려될 때 렙톤 질량이 \(Q=2/3\)을 만족하도록 강제한다. 이 결과는 고이데 관계식을 질량 계층, 우주 상수, 그리고 기하학적 항등식 \(S_{\mathrm{odd}}\varphi^2 \approx \pi\)와 통일시키며, 이들 모두는 YUCT의 단일한 대수적 폐루프로부터 흘러나온다.
	
	% ----------------------------------------------------------------
	
	\section{반증 가능성과 미래 검증}
	
	이 부록의 예측들은 구체적이고 반증 가능하다:
	\begin{itemize}
		\item 방정식 \eqref{eq:hierarchy}은 페르미온 자유도의 수를 바꾸는 표준 모형의 어떤 확장도 약전자기 척도를 이동시킬 것임을 함의한다. \(m_W\)와 \(M_{\mathrm{Pl}}\)(\(G\)를 통해)의 미래 정밀 측정이 이 관계를 검증할 수 있다.
		\item 방정식 \eqref{eq:lambda-suppression}은 \(\rho_\Lambda\)를 지평선 면적을 통해 \(H_0\)와 묶는다. \(H_0\)가 더 정확히 측정됨에 따라(예: Euclid, Roman, 또는 CMB‑S4), 예측된 \(\Omega_\Lambda\)는 관측값과 일치해야 한다. 유의미한 이탈은 모형을 반증할 것이다.
		\item 중입자 계층 \eqref{eq:baryon-hierarchy}는 \(H_0\), \(\Omega_b\), 그리고 기본 질량들의 정밀도를 향상시킴으로써 직접적으로 검증 가능하다. 현재의 불확실성은 이미 일치를 설득력 있게 만들며; 미래 데이터는 그것을 확증하거나 배제할 것이다.
	\end{itemize}
	이 예측들을 조정할 수 있는 자유 매개변수는 없다——그것들은 YUCT의 대수적 폐루프에 의해 고정되어 있다.
	
	\subsection{예측: \(A=298\)에서의 안정성의 섬}
	
	YUCT의 핵 질량에 대한 조정 깊이 분석은 알려진 이중 마법 핵 \(^{208}\mathrm{Pb}\)과 예상되는 다음 껍질 닫힘 사이의 \(L_{\mathrm{eff}}\) 공명 네트워크에 상당한 틈이 있음을 드러낸다. 안정된 핵 질량이 기본 비 \(3/2\)의 거듭제곱과 정렬해야 한다는 요구는 구체적인 예측으로 이끈다: 핵 안정성의 국소적 최대값이 질량수 \(A = 298 \pm 2\)에서 발생해야 하며, 이는 원소 \(Z \approx 120\)--\(122\)에 해당한다. 이 예측은 상세한 껍질 모형 퍼텐셜과 무관하며, 조정 깊이의 이산적 대수 구조에만 의존한다. 수명이 \(1\ \mu\text{s}\)를 초과하는 그러한 핵의 합성은 질량의 조정 기반 기원에 대한 강력한 증거를 제공할 것이다.
	
	\section{결론}
	
	우리는 야쿠셰프 통합 조정 이론이 계층 문제와 우주 상수 문제 양자 모두에 대해 매개변수 없이 정량적인 해결책을 제공함을 보였다. \(10^{17}\)과 \(10^{122}\)라는 거대한 비율은 각각 기본 페르미온의 수와 우주 지평선의 엔트로피로 거슬러 올라가며, 보편 상수 \(S_{\mathrm{odd}}\)와 \(S_{\mathrm{even}}\)을 통해 연결된다. 중입자 질량 계층은 두 척도 사이의 통일하는 다리 역할을 한다. 모든 예측은 다가오는 관측 데이터로 반증 가능하다. 따라서 YUCT는 현대 물리학의 가장 큰 두 미세 조정 퍼즐에 대한 우아하고 검증 가능한 해결책을 제공한다.
	
	\section*{결론 및 전망}
	
	마지막으로, YUCT 대수적 폐루프의 내적 정합성은 고정밀 기하학적 우연에 의해 놀랍게도 예시된다. 페르미온 질량 계층을 지배하는 상수들 \(S_{\mathrm{odd}}=6/5\)와 황금비 \(\varphi\)는 함께 원주율 \(\pi\)를 단지 \(1.5\times10^{-5}\)의 오차로 근사한다:
	\[
	S_{\mathrm{odd}}\varphi^2 \approx \pi .
	\]
	이는 고전적 유리수 근사 \(22/7\)보다 한 자릿수 더 정확하다. 페르미온 질량의 반정수 양자화와 함께, 이것은 동일한 이산적 대수 구조가 입자 질량의 이산 스펙트럼과, 조정된 물질로부터의 연속 기하학의 창발(부록 Ehrenfest) 모두를 지배한다는 것을 보여준다. 이러한 이중 우연은 우연적이기 어려우며 물리 법칙의 조정적 기원에 대한 설득력 있는 증거를 제공한다.
	
	\section*{감사의 말}
	
	저자는 YUCT 세미나 참가자들의 자극적인 토론에 감사한다.
	
	\section*{데이터 이용 가능성}
	
	모든 유도와 보조 계산은 YPSDC 리포지토리 \url{https://github.com/Alexey-Yakushev-YUCT/YPSDC}에서 이용 가능하다.
	
	\begin{thebibliography}{99}
		\bibitem{AppendixL} A.V. Yakushev, ``Appendix L: Fractal Coordination Error Scaling — A Universal Law from DNA to Cosmology,'' in \emph{YUCT}, Zenodo (2026).
		\bibitem{AppendixY} A.V. Yakushev, ``Appendix Y: Mathematical Foundations of YUCT — A Derivation from First Principles,'' in \emph{YUCT}, Zenodo (2026).
		\bibitem{AppendixM} A.V. Yakushev, ``Appendix M: YUCT Cosmology — Inflation as a Coordination Phase Transition,'' in \emph{YUCT}, Zenodo (2026).
		\bibitem{AppendixΩ} A.V. Yakushev, ``Appendix $\Omega$: The Global Rotation of the Universe and Its New Minimum Age from the Dipole Turning Radius,'' in \emph{YUCT}, Zenodo (2026).
		\bibitem{AppendixFerra1.2} A.V. Yakushev, ``Appendix Ferra1.2: Universal Coordination Constants for Ordered and Disordered Phases,'' in \emph{YUCT}, Zenodo (2026).
		\bibitem{Koide1983} Y.~Koide, ``A New View of Quark and Lepton Mass Hierarchy,''
		\emph{Phys. Rev. D} \textbf{28}, 252--254 (1983).
		\bibitem{Koide1982} Y.~Koide, ``Fermion‑Boson Two‑Body Model of Quark and Lepton Masses,''
		\emph{Lett. Nuovo Cimento} \textbf{34}, 201--205 (1982).
		\bibitem{Foot1994} R.~Foot, ``A Note on the Koide Lepton Mass Relation,''
		\emph{Phys. Lett. B} \textbf{326}, 307--311 (1994).
	\end{thebibliography}
	
\end{document}